% Generated draft: requires mathematical and source review.
\documentclass[11pt]{article}
\usepackage[margin=25mm]{geometry}
\usepackage{fontspec}
\setmainfont{Latin Modern Roman}
\usepackage{amsmath,amssymb,mathtools,unicode-math}
\setmathfont{Latin Modern Math}
\usepackage{enumitem,needspace,xurl,hyperref}
\setlength{\parindent}{0pt}
\setlength{\parskip}{5pt}
\setlength{\emergencystretch}{3em}
\newcommand{\N}{\mathbb{N}}
\newcommand{\Z}{\mathbb{Z}}
\newcommand{\Q}{\mathbb{Q}}
\newcommand{\R}{\mathbb{R}}
\newcommand{\C}{\mathbb{C}}
\newcommand{\F}{\mathbb{F}}
\newcommand{\A}{\mathbb{A}}
\newcommand{\PP}{\mathbb{P}}
\newcommand{\E}{\mathbb{E}}
\newcommand{\eps}{\varepsilon}
\newcommand{\dd}{\,\mathrm{d}}
\begin{document}
\section*{UnsolvedMath ID: 30001503}
\textbf{Source problem code:} catalogue entry
\subsection{Definitions and mathematical statement}
Let \(P\) be a circle packing in \(\C\), described in the supplied statement as locally finite. Under the usual Euclidean convention, a circle \(C\) has radius \(r(C)>0\) and unsigned curvature \(\operatorname{curv}(C)=1/r(C)\); a circle packing consists of circles whose bounded open disks have pairwise disjoint interiors. These conventions require confirmation against the source, which does not explicitly define them. For a bounded Borel set \(E\subseteq\C\) and \(T>0\), write
\[
N_P(T,E)=\#\{C\in P:C\cap E\ne\varnothing,\ \operatorname{curv}(C)<T\}.
\]
Here intersection refers to the circle itself. Whether the intended source convention instead counts intersection with the enclosed disk remains unresolved.

The supplied question asks whether, for each such packing \(P\), there exist a Borel measure \(\omega_P\) on \(\C\) and an exponent \(\alpha\), both independent of \(E\) and \(T\), such that for every suitable bounded Borel set \(E\),
\[
N_P(T,E)\sim\omega_P(E)T^\alpha\qquad(T\to\infty).
\]
It also asks what analogous counting and distribution results hold for circle packings beyond Apollonian packings.

The supplied record does not define ``suitable,'' specify the allowed exponents, impose nontriviality or local finiteness conditions on \(\omega_P\), or explain how to interpret the displayed asymptotic when \(\omega_P(E)=0\) or \(\omega_P(E)=\infty\). For \(0<\omega_P(E)<\infty\), the conventional meaning is
\[
\lim_{T\to\infty}\frac{N_P(T,E)}{\omega_P(E)T^\alpha}=1.
\]
A potentially useful formulation covering zero measure sets would instead require \(T^{-\alpha}N_P(T,E)\to\omega_P(E)\), but the record does not establish that this is the intended formulation. No boundary regularity condition, such as \(\omega_P(\partial E)=0\), is supplied and none is imposed here.

The meaning of local finiteness is particularly consequential: finiteness of all circles meeting each compact set differs from finiteness only after a curvature cutoff. The record does not resolve this distinction. Its background also describes an unrestricted class of countable packings, whereas its displayed question assumes local finiteness. Thus this is a statement-only draft of the source question with unresolved hypotheses, rather than a fully specified universal assertion.
\subsection{Short English statement}
The question asks whether each locally finite planar circle packing admits a common power law for counting circles of curvature less than \(T\) that meet a suitable bounded region, with a Borel measure describing spatial distribution. It also asks for extensions beyond Apollonian packings. The intended local finiteness condition, admissible regions, measure conditions, and treatment of zero measure regions need clarification. The supplied dated curation labels the problem partially solved; this draft does not independently verify that assessment.
\subsection{Sources}
\begin{itemize}
\item Supplied local UnsolvedMath registry record, permanent ID \(30001503\), title ``Asymptotic Counting in General Circle Packings,'' source problem code \texttt{catalogue entry}, in the Oberwolfach Reports Open Problems collection. The proposed year is recorded as 2010; no individual proposer is supplied.
\item Supplied report citation: \emph{Homogeneous Dynamics and Number Theory (2011)}, Oberwolfach Reports, DOI 10.4171/owr/2010/29, pp.\ 1703--1759. \url{https://doi.org/10.4171/owr/2010/29}. The citation's 2011 designation and the 2010 designation in the DOI and proposed-year field are retained without reconciliation.
\item Supplied literature source labeled ``Kontorovich Oh'': \url{https://arxiv.org/abs/0811.2236}.
\item Supplied literature source labeled ``circle packing survey'': \url{https://arxiv.org/abs/1312.1383}.
\item Literature assessment dated exactly \texttt{2026-08-21}: the supplied curation labels the problem partially solved and states that counting and distribution are proved for Apollonian packings and broad classes of geometrically finite Kleinian invariant packings, while the unrestricted class of countable circle packings is not classified. This is supplied dated curation, not an independently verified assessment; its time-sensitive status and cited claims require independent verification.
\item The registry labels the statement \texttt{corrected\_verified} and reports that its mathematical content was checked against the official Oberwolfach report, with notation, context, and question form normalized. This is a supplied verification claim; no source was visited or independently checked for this draft.
\end{itemize}
\end{document}
